% !PRINTED 316
\weberpage{316}
Die beiden Substitutionen $1, \gamma_0$ bilden eine Gruppe $Q_0$. Irgend eine Substitution $\gamma_1$ von $Q$, die nicht in $Q_0$ vorkommt, führt 0, 1 in zwei davon verschiedene Elemente, etwa 2, 3, über, und durch die zwei Substitutionen $Q_0\gamma_1$ geht 0, 1 in 2, 3 oder in 3, 2 über. Wenn ferner durch $\gamma_2$ das Paar 0, 1 in ein drittes Paar 4, 5 übergeht, so ist 4, 5 sowohl von 0, 1 als von 2, 3 verschieden, und so können wir $Q$ in die Nebengruppen zerlegen:

\[
Q = Q_0 + Q_0\gamma_1 + Q_0\gamma_2 + Q_0\gamma_3,
\]

wodurch die acht Elemente (1) in vier Paare

\weq{2}{0, 1;\quad 2, 3;\quad 4, 5;\quad 6, 7}

zerlegt sind, so dass durch die beiden Substitutionen einer der Nebengruppen $Q_0\gamma_i$ das Paar 0, 1 in eines dieser vier Paare übergeht.

Ist nun $k$ irgend eine Substitution aus der Gruppe

\[
P = T_{24}\,Q,
\]

so lassen sich die Elemente $\tau_0, \tau_1, \tau_2, \tau_3$ aus $T_{24}$ und $\gamma_0', \gamma_1', \gamma_2', \gamma_3'$ aus $Q$ so bestimmen, dass

\[
k = \tau_0\gamma_0', \quad \gamma_1 k = \tau_1\gamma_1', \quad \gamma_2 k = \tau_2\gamma_2', \quad \gamma_3 k = \tau_3\gamma_3'
\]

oder

\[
k = \tau_0\gamma_0' = \gamma_1^{-1}\tau_1\gamma_1' = \gamma_2^{-1}\tau_2\gamma_2' = \gamma_3^{-1}\tau_3\gamma_3',
\]

und diese Darstellung zeigt, dass durch $k$ irgend eines der Paare (2) in ein Paar desselben Systems übergeführt wird; denn z.\,B. durch $\gamma_1^{-1}$ wird 2, 3 in 0, 1 übergeführt, durch $\tau_1$ bleibt 0, 1 ungeändert und durch $\gamma_1'$, was zu $Q$ gehört, wird 0, 1 in eines der Paare (2) übergeführt.

Die Gruppe $T_{24}\,Q$ ist also imprimitiv, und zwar so, dass wir vier Systeme der Imprimitivität haben. Die Wurzeln 0, 1 genügen einer quadratischen Gleichung, deren Coëfficienten von den Wurzeln einer biquadratischen Gleichung abhängen.

Gehen wir nun zu dem Falle 2) über, in dem die Elemente je eines Tripels 1, 2, 3 und eines Quadrupels 4, 5, 6, 7 durch $T_{24}$ transitiv unter einander verbunden sind. Die acht Grössen (1) zerfallen hier in zwei Quadrupel:

\weq{3}{0, 1, 2, 3;\quad 4, 5, 6, 7,}

und durch die Substitutionen von $Q$ werden die Grössen eines dieser Quadrupel entweder nur unter sich vertauscht, oder sie werden in die Grössen des anderen Quadrupels übergeführt.
